% ------------------------------------------------------------------------------
% Introdução
% ------------------------------------------------------------------------------

\chapter{Introdução}
\label{chap_introducao}

O passeio aleatório simples em uma dimensão (\textit{random walk} 1D) e um dos modelos mais elementares e, ao mesmo tempo, mais relevantes da modelagem estocastica. Mesmo com regras locais simples, esse processo reproduz comportamentos globais importantes em fisica estatistica, difusao, teoria de filas e simulacao computacional \cite{mathworld_random_walk_1d}.

Nesta pratica, o random walk foi implementado a partir de um gerador linear congruencial (LCG), e nao por rotinas prontas de pseudoaleatoriedade, para reforcar o entendimento sobre a geracao de variaveis aleatorias e seu impacto em simulacoes de Monte Carlo. A proposta do trabalho inclui cinco etapas: (i) construcao do gerador e da caminhada, (ii) simulacao de 10 trajetorias com 10000 passos, (iii) analise do desvio quadratico medio em escala log-log com ajuste de lei de potencia, (iv) teste de hipotese para o expoente ajustado e (v) verificacao do Teorema Central do Limite (TCL) com 50000 caminhadas adicionais.

O objetivo geral e verificar se os resultados numericos produzidos pela implementacao sao compativeis com o comportamento teorico esperado de um random walk simples sem vies. Em particular, espera-se que o desvio quadratico medio obedeça aproximadamente a relacao $\langle R^2(t)\rangle \propto t^{\alpha}$ com $\alpha \approx 1$ e que os deslocamentos finais normalizados apresentem comportamento aproximadamente gaussiano para grande numero de passos.

Este relatorio esta organizado da seguinte forma: o Capitulo~\ref{chap_fundamentacao_teorica} apresenta os conceitos teoricos essenciais; o Capitulo~\ref{chap_metodologia} descreve os procedimentos computacionais adotados; o capitulo de resultados apresenta as evidencias numericas e graficas; e, por fim, o Capitulo~\ref{chap_conclusao} sintetiza as principais conclusoes e desdobramentos.

